\scanpage{42}
Enfin, on devrait donner un
\uline{Th. 5} Subdivision barycentrique de $\psi_n(C_n)$, et description par ens. pseudo-simplicial \ldots.

\uline{NB} La donnée d'un homéom. $X_n\simeq\psi_n\varphi_n(X_n)$ revient à la donnée, pour tt $i$ et tte composante connexe $\widetilde X_i^\alpha$ de $\widetilde X_i$ (``$i$-cellule'') d'un homéomorphisme
\[ \widetilde X_i^\alpha\xrightarrow{\sim}\mathrm{C\hat{o}ne}\,(\partial\widetilde X_i^\alpha) \]
induisant l'identité sur le bord.

La subdivision barycentrique ne fait pas tjrs une triangulation.
On pourrait expliciter la subdivision barycentrique itérée $2^{\text{\unsure{ième|ièm.}}}$, qui correspond, elle, à une triangulation. Il faudrait donc expliciter \deleted{la} le squelette simplicial (ensemble $+$ ens de parties finies, stable par inclusion) de cet espace triangulé, en termes de la cascade $C_n$, et donner une description sur ce \deleted{\unsure{s}} squelette de la filtration qu'on a sur sa réalisation topologique. Cel fera bien encore un th. 6 (et peut-être un th. 7\ldots)
% archive sheet number: 41
